%%%%%%%%%%%%%%%%%%%%%%%%%%%%%%%%%%%%%%%%%%%%%%%%%%%%%%%%%%%%%%%%%%%%%%%
%% ARCHIVO: styles/tikz_styles.tex
%% Estilos para diagramas de flujo con TikZ
%%%%%%%%%%%%%%%%%%%%%%%%%%%%%%%%%%%%%%%%%%%%%%%%%%%%%%%%%%%%%%%%%%%%%%%

\usepackage{tikz}
\usetikzlibrary{
    shapes.geometric,
    arrows.meta,
    calc,
    positioning,
    shadows.blur,
    fit,
}

% === DEFINICIÓN DE LA FORMA 'document' ===
\makeatletter
\pgfdeclareshape{document}{
    \inheritsavedanchors[from=rectangle]
    \inheritanchorborder[from=rectangle]
    \inheritanchor[from=rectangle]{center}
    \inheritanchor[from=rectangle]{north}
    \inheritanchor[from=rectangle]{south}
    \inheritanchor[from=rectangle]{west}
    \inheritanchor[from=rectangle]{east}
    \backgroundpath{
        \southwest \pgf@xa=\pgf@x \pgf@ya=\pgf@y
        \northeast \pgf@xb=\pgf@x \pgf@yb=\pgf@y
        \pgf@xc=\pgf@xb \advance\pgf@xc by-10pt
        \pgf@yc=\pgf@yb \advance\pgf@yc by-10pt
        \pgfpathmoveto{\pgfpoint{\pgf@xa}{\pgf@ya}}
        \pgfpathlineto{\pgfpoint{\pgf@xa}{\pgf@yb}}
        \pgfpathlineto{\pgfpoint{\pgf@xc}{\pgf@yb}}
        \pgfpathlineto{\pgfpoint{\pgf@xb}{\pgf@yc}}
        \pgfpathlineto{\pgfpoint{\pgf@xb}{\pgf@ya}}
        \pgfpathclose
        \pgfpathmoveto{\pgfpoint{\pgf@xc}{\pgf@yb}}
        \pgfpathlineto{\pgfpoint{\pgf@xc}{\pgf@yc}}
        \pgfpathlineto{\pgfpoint{\pgf@xb}{\pgf@yc}}
    }
}
\makeatother

% === ESTILOS DE NODOS ===
\tikzset{
    startstop/.style={
        ellipse, 
        draw=black!80, 
        fill=gray!25,
        blur shadow={shadow blur steps=5},
        minimum width=3.5cm, 
        minimum height=1.0cm, 
        text centered, 
        font=\bfseries,
        line width=1.2pt
    },
    io/.style={
        trapezium, 
        trapezium left angle=70, 
        trapezium right angle=110,
        draw=blue!70, 
        fill=blue!20,
        blur shadow={shadow blur steps=5},
        minimum width=3cm, 
        minimum height=0.9cm, 
        text centered,
        font=\small\sffamily,
        line width=1pt
    },
    process/.style={
        rectangle, 
        draw=orange!70, 
        fill=orange!20,
        blur shadow={shadow blur steps=5},
        minimum width=2cm, 
        minimum height=1.0cm, 
        text centered,
        font=\small\sffamily,
        rounded corners=3pt, 
        line width=1pt
    },
    decision/.style={
        diamond, 
        draw=green!60!black, 
        fill=green!20,
        blur shadow={shadow blur steps=5},
        minimum width=1.8cm, 
        minimum height=0.9cm, 
        text centered,
        font=\small\sffamily,
        aspect=2.5, 
        inner sep=0pt,
        line width=1pt
    },
    doc/.style={
        document, 
        draw=yellow!70!orange, 
        fill=yellow!25,
        blur shadow={shadow blur steps=5},
        minimum width=3.2cm, 
        minimum height=0.9cm, 
        text centered,
        font=\small\sffamily,
        line width=1pt
    },
    error/.style={
        document,
        draw=red!70,
        fill=red!15,
        blur shadow={shadow blur steps=5},
        minimum width=3.2cm,
        minimum height=0.9cm,
        text centered,
        font=\small\sffamily\bfseries,
        line width=1.2pt
    },
    arrow/.style={
        -Stealth,
        thick,
        draw=black!70,
        line width=1pt
    },
    arrow_yes/.style={
        arrow,
        draw=green!60!black
    },
    arrow_no/.style={
        arrow,
        draw=red!60!black
    },
    label_node/.style={
        font=\scriptsize\sffamily\bfseries,
        fill=white,
        inner sep=2pt
    },
    module_box/.style={
        rectangle,
        draw=purple!70,
        thick,
        dashed,
        rounded corners=5pt,
        inner sep=10pt
    }
}

% === DEFINICIONES DE TEXTO ===
\def\isTrue{Sí}
\def\isFalse{No}